\section{Functions} 

\subsection{Continuous Functions}

	Recall what a \href{def:continuous_func}{continuous function} is. 

	\begin{lemma}[Properties of Continuous Functions]
		Continuous functions satisfy the following properties. 
		\begin{enumerate}
			\item \textit{Identity}. $\id: (X, \T_1) \rightarrow (X, \T_2)$ is continuous if $\T_1 \supset \T_2$. 
			\item \textit{Constant}. A constant function $c_{y_0}: (X, \T) \rightarrow (Y, \T_2)$ defined $c(x) = y_0$ is always continuous, regardless of the topologies. 
		  \item \textit{Composition}. If $f: X \rightarrow Y$ and $g: Y \rightarrow Z$ is continuous, then $g \circ f :X \rightarrow Z$ is continuous. 
			\item \textit{Restriction}. 
			\item \textit{Injection}. 
			\item \textit{Product}. 
		\end{enumerate}
	\end{lemma}
  \begin{proof}
    Listed. 
    \begin{enumerate}
			\item \textit{Identity}. If we take an open set $U \in \T_2$, its preimage is the same set $U$, which is guaranteed to be in $\T_1$ since $\T_1$ is finer than $\T_2$. 

			\item \textit{Constant}. Trivial. 

			\item \textit{Composition}. Let $U \subset Z$ be open. Then $g^{-1}(Z)$ is open in $Y$, and so $f^{-1}(g^{-1}(Z))$ is open in $X$. 

			\item \textit{Restriction}. Let us take an open set $U$ in $Y$. Then it is of the form $V \cap Y$ for some $V$ open in $X$. Therefore taking the preimage gives 
				\begin{equation}
					f|_{Z}^{-1} (U) = f^{-1} (U) = f^{-1} (V \cap Y) = f^{-1} (V) \cap f^{-1} (Y) = f^{-1} (V) \cap Z
				\end{equation}
				where $f^{-1} (V)$ is open by continuity of $f$, and so the intersection is open. 

			\item \textit{Injection}. 
    \end{enumerate}
  \end{proof}

  \begin{theorem}[Arithmetic of Real Continuous Functions]
		\label{thm:arithmetic_continuous}
    If $X$ is a topological space, and if $f, g: X \longrightarrow \mathbb{R}$ are continuous at $x \in X$, then 
		\begin{enumerate}
		  \item $f + g$ is continuous at $x$
			\item $cf$ is continuous at $x$ for all $c \in \mathbb{R}$. 
			\item $fg$ is continuous at $x$. 
			\item $f/g$ is continuous at $x$ if $g \neq 0$ for some open neighborhood $U_x \ni x$. 
		\end{enumerate}
  \end{theorem}
	\begin{proof}
	  
	\end{proof}

  Given these results, one may wonder whether---just like how we restricted a continuous function to a smaller continuous function---we can ``extend'' a function to a larger function. However, this is not always true. 

  \begin{example}[Combining Continuous Functions May not be Continuous]
    Let us take $\mathbb{R}$ and divide it into $\mathbb{Q}$ and $(\mathbb{R} \setminus \mathbb{Q}) \setminus \{0\}$. Then let us define 
    \begin{align}
      f: \mathbb{Q} \rightarrow \mathbb{R} & f(x) = 0 \\ 
      g: (\mathbb{R} \setminus \mathbb{Q}) \setminus \{0\} \rightarrow \mathbb{R} & g(x) = x
    \end{align}
    Then $f$ and $g$ are trivially continuous, but taking the function 
    \begin{equation}
      h(x) \coloneqq \begin{cases} 
        f(x) = 0 & \text{ if } x \in \mathbb{Q} \\ 
        g(x) = x & \text{ if } x \not\in \mathbb{Q}
      \end{cases}
    \end{equation}
    which is not continuous.\footnote{Inspired from \href{https://math.stackexchange.com/questions/4034361/any-counter-example-of-pasting-lemma}{here}. }
  \end{example} 

  But not all hope is lost. It does turn out that under certain conditions, we can in fact construct such continuous functions. 

  \begin{lemma}[Pasting Lemma, Gluing Lemma]
		\label{thm:pasting_lemma}
    Let $X = A \cup B$, where $A, B$ are closed in $X$. Let $f: A \longrightarrow Y$ and $g: B \longrightarrow Y$ be continuous. If 
    \begin{equation}
      f(x) = g(x) \text{ for all } x \in A \cap B
    \end{equation}
    Then $f$ and $g$ can be combined to form a continuous function $h: X \longrightarrow Y$, defined
    \begin{equation}
      h(x) \coloneqq \begin{cases}
        f(x) & x \in A \setminus B \\
        f(x) \text{ or } g(x) & x \in A \cap B \\
        g(x) & x \in B \setminus A
      \end{cases}
    \end{equation}

    \begin{figure}[H]
      \centering 
      \begin{tikzpicture}[
        scale=1.2,
        set/.style={draw, thick},
        mapping/.style={->, thick},
        dashed mapping/.style={->, thick, dashed}
      ]
        % Set A with red north west lines pattern (smooth, asymmetric blob shape, vertically elongated)
        \begin{scope}
            \clip (-0.5,1.8) .. controls (-1.2,2.25) and (-2,1.8) .. 
                  (-2.3,0.6) .. controls (-2.5,-0.45) and (-1.8,-1.35) .. 
                  (-1,-1.05) .. controls (-0.2,-0.75) and (0.8,-0.3) .. 
                  (1.5,0.45) .. controls (2,1.2) and (1.5,2.1) .. 
                  (0.7,2.25) .. controls (0,2.3) and (-0.5,1.8) .. 
                  (-0.5,1.8) -- cycle;
            \fill[pattern=north west lines, pattern color=red!70] (-3,-3) rectangle (3,3);
        \end{scope}
        \draw[set] (-0.5,1.8) .. controls (-1.2,2.25) and (-2,1.8) .. 
                  (-2.3,0.6) .. controls (-2.5,-0.45) and (-1.8,-1.35) .. 
                  (-1,-1.05) .. controls (-0.2,-0.75) and (0.8,-0.3) .. 
                  (1.5,0.45) .. controls (2,1.2) and (1.5,2.1) .. 
                  (0.7,2.25) .. controls (0,2.3) and (-0.5,1.8) .. 
                  (-0.5,1.8) -- cycle;
        \node at (-1,0.75) {$A$};
        
        % Set B with blue north east lines pattern (asymmetric blob shape, vertically elongated)
        \begin{scope}
            \clip (-0.5,-1.95) .. controls (-1.8,-2.25) and (-2.5,-1.2) .. 
                  (-1.6,0) .. controls (-0.8,0.9) and (0.2,1.2) .. 
                  (1,0.75) .. controls (1.8,0.3) and (2.2,-0.45) .. 
                  (1.7,-1.5) .. controls (1,-2.25) and (0.3,-1.8) .. 
                  (-0.5,-1.95) -- cycle;
            \fill[pattern=north east lines, pattern color=blue!70] (-3,-3) rectangle (3,3);
        \end{scope}
        \draw[set] (-0.5,-1.95) .. controls (-1.8,-2.25) and (-2.5,-1.2) .. 
                  (-1.6,0) .. controls (-0.8,0.9) and (0.2,1.2) .. 
                  (1,0.75) .. controls (1.8,0.3) and (2.2,-0.45) .. 
                  (1.7,-1.5) .. controls (1,-2.25) and (0.3,-1.8) .. 
                  (-0.5,-1.95) -- cycle;
        \node at (1.2,-1.2) {$B$};
        
        % Add the intersection label
        \node at (0,0) {$A \cap B$};
        
        % Draw |___| shaped arrow at the bottom
        \draw[mapping] (0,-2.5) -- (0,-3) -- (5.3,-3) -- (5.3,-2.5) node[midway, left] {$h$};
        
        % Rectangle around right side (set Y)
        \draw[set] (3.8,-2) rectangle (6.8,2);
        \node at (6.4,-1.7) {$Y$};
        
        % Set f(A) on right with red north west lines pattern
        \begin{scope}
            \clip (5,1.5) .. controls (4.3,1.8) and (4,1.2) .. 
                  (4.2,0.5) .. controls (4.5,-0.3) and (5,-0.6) .. 
                  (5.8,-0.2) .. controls (6.3,0.2) and (6.5,0.8) .. 
                  (6.2,1.3) .. controls (5.8,1.7) and (5.5,1.3) .. 
                  (5,1.5) -- cycle;
            \fill[pattern=north west lines, pattern color=red!70] (3.5,-1) rectangle (7,2);
        \end{scope}
        \draw[set] (5,1.5) .. controls (4.3,1.8) and (4,1.2) .. 
                  (4.2,0.5) .. controls (4.5,-0.3) and (5,-0.6) .. 
                  (5.8,-0.2) .. controls (6.3,0.2) and (6.5,0.8) .. 
                  (6.2,1.3) .. controls (5.8,1.7) and (5.5,1.3) .. 
                  (5,1.5) -- cycle;
        \node at (5.3,0.5) {$f(A)$};
        
        % Set g(B) on right with blue north east lines pattern
        \begin{scope}
            \clip (5,-1.2) .. controls (4.2,-1.5) and (4,-0.8) .. 
                  (4.3,-0.3) .. controls (4.7,0.3) and (5.3,0.5) .. 
                  (5.9,0.1) .. controls (6.5,-0.3) and (6.4,-1) .. 
                  (5.8,-1.4) .. controls (5.4,-1.6) and (5,-1) .. 
                  (5,-1.2) -- cycle;
            \fill[pattern=north east lines, pattern color=blue!70] (3.5,-2) rectangle (7,1);
        \end{scope}
        \draw[set] (5,-1.2) .. controls (4.2,-1.5) and (4,-0.8) .. 
                  (4.3,-0.3) .. controls (4.7,0.3) and (5.3,0.5) .. 
                  (5.9,0.1) .. controls (6.5,-0.3) and (6.4,-1) .. 
                  (5.8,-1.4) .. controls (5.4,-1.6) and (5,-1) .. 
                  (5,-1.2) -- cycle;
        \node at (5.3,-0.8) {$g(B)$};
        
        % Dotted arrows for direct mappings
        \draw[dashed mapping] (-0.5,1) to[out=20, in=160] node[midway, above] {$f$} (4.3,1);
        \draw[dashed mapping] (-0.3,-1.5) to[out=-20, in=200] node[midway, below] {$g$} (4.3,-1);
        \draw[dashed mapping] (0,0) to[out=0, in=180] node[midway, above] {$f$ or $g$} (4.8,0);
      \end{tikzpicture}
      \caption{Visual of the pasting lemma.} 
      \label{fig:gluing_lemma}
    \end{figure}
  \end{lemma}

\subsection{Open and Closed Maps}

  Open and closed functions map open/closed sets to open/closed sets, unlike continuous functions which take the preimage. However, they do are not natural and most maps are not open nor closed, so this is a pretty special condition. 

  \begin{definition}[Open, Closed Maps]
    A map $f: X \rightarrow Y$ is said to be 
    \begin{enumerate}
      \item \textbf{open} if it maps open sets of $X$ to open sets of $Y$. 
      \item \textbf{closed} if it maps open sets of $X$ to closed sets of $Y$. 
    \end{enumerate}
    Note that open and closed maps are completely independent. A map may be open, closed, neither, or both. 
  \end{definition}

  \begin{example}[Open but Not Closed]
    The projection $\pi_1: X \times Y \rightarrow X$ is an open map but but closed. Consider $\pi_1: \mathbb{R}^2 \rightarrow \mathbb{R}$ with $S = \{(x, y) \in \mathbb{R}^2 \mid xy = 1 \}$. Then $\pi_1 (S) = \mathbb{R} \setminus \{0\}$, which is not closed.\footnote{In open maps, the typical behavior is that points are ``copied,'' i.e. for projections, the preimage of $\pi_1^{-1} (x) = x \times Y$, where all $y \in Y$ are copied.}
  \end{example}

  \begin{example}[Closed but Not Open]
    $f: \mathbb{R} \rightarrow \mathbb{R}$ with $f(x) = x^2$ is closed but not open since $f(\mathbb{R}) = [0, +\infty)$ which is not open. 
  \end{example}

  Note that an open map or a closed map (with continuous and surjective) are trivially quotient maps. Since given a $U \subset Y$ with $f^{-1} (U)$ open, then by definition $U = f(f^{-1} (U))$ is open by definition. 

  \begin{theorem}[Open/Closed Maps are Quotient Maps]
    If $p: X \rightarrow Y$ is a surjective, continuous map that is either open or closed (that is, maps open sets to open sets or closed sets to closed sets), then $p$ is a quotient map.\footnote{Note however, that the converse is not true; there exists quotient maps that are neither open nor closed. }
  \end{theorem} 

  \begin{example}[Quotient Maps that are Neither Open Nor Closed]
    
  \end{example}

\subsection{Homeomorphisms}

	Recall what a \href{def:homeomorphism}{homeomorphism} is. Not only does a homeomorphism give a bijective correspondence between points in $X$ and $Y$, but it also determines a bijection between their topologies! This bijection then allows two spaces that are homeormophic to have the same topological properties, such as separability, countability, compactness, and connectedness. 

	\begin{lemma}[Properties of Homeomorphisms]
	  \begin{enumerate}
	    \item \textit{Identity}. $\id_X: (X, \T) \to (X, \T)$ is always a homeomorphism. 
			\item \textit{Composition}. If $f: X \to Y$ and $g: Y \to Z$ are homeomorphisms, then $g \circ f: X \to Z$ is a homeomorphism. 
			\item \textit{Restriction}. If $f: X \rightarrow Y$ is a homeomorphism, then for any $x \in X$, the restriction 
				\begin{equation}
					f|_{X \setminus \{x\}} : X \setminus \{x\} \rightarrow Y \setminus \{f(x)\}
				\end{equation}
				is also a homeomorphism. 
	  \end{enumerate}
	\end{lemma}
	\begin{proof}
	  
	\end{proof}

  \begin{example}[Comparability and Homeomorphic Spaces]
    Consider the set $X = \{a, b\}$ with the two topologies $\T_3 = \{\emptyset, \{a\}, X\}$ and $\T_4 = \{\emptyset, \{b\}, X\}$. They are not comparable but they seem ``similar'' in a way in that if we swap all the $a$'s and $b$'s in $\T_3$, then we get $\T_4$. We can make this rigorous by defining $f: (X, \T_3) \rightarrow (X, \T_4)$ with $f(a) = b, f(b) = a$, and showing that it is a homeomorphism. 
  \end{example}

	Note that we may have functions that are bijective and continuous, but not bicontinuous. In order to construct such examples one of the easiest things we can do is just endow the codomain with the discrete topology, which guarantees continuity. 

  \begin{example}[Bijective and Continuous but not Homeomorphism]
    $\mathbb{Z}$ and $\mathbb{Q}$ are countable sets, so there is a bijection between them. If we give each of them the metric topology, $\mathbb{Z}$ ends up having the discrete topology (take the $0.5$-ball around each integer), whereas for $\mathbb{Q}$, we will see later that by the density of the rationals there are an infinite number of rationals in $(q - r, q + r)$ for $q \in \mathbb{Q}$. Note that this bijection $f: \mathbb{Z} \rightarrow \mathbb{Q}$ is continuous (since $\mathbb{Z}$ has discrete topology) but not bicontinuous. 
  \end{example}


	Let's talk about this in categorical notions. The collection of all topological spaces, together with continuous maps as morphisms, forms the category of topological spaces $\mathbf{Top}$. The isomorphisms in $\mathbf{Top}$ are precisely the homeomorphisms. Therefore, being homeomorphic is an equivalence relation on the objects of $\mathbf{Top}$, and its equivalence classes are called homeomorphism classes, or equivalently, isomorphism classes in $\mathbf{Top}$. Since the collection of all topological spaces is a proper class, $\mathbf{Top}$ is a large category. Nevertheless, for any two spaces (X,Y), the collection $\operatorname{Hom}_{\mathbf{Top}}(X,Y)$ of continuous maps is a set, making $\mathbf{Top}$ locally small.

	\begin{definition}[Embedding]
	  A function $f: X \to Y$ is an \textbf{embedding} if $f$ is a homeomorphism onto its image. That is, $f: X \to f(X)$ is a homeomorphism. 
	\end{definition}

\subsection{Local Homeomorphisms} 

	\begin{definition}[Locally Homeomorphic]
	  
	\end{definition}

  \begin{example}[Homeomorphism Classes of 2D Manifolds]
    There is an infinite family of 2-dimensional manifolds, call them $M$ and $N$, and each set in each family is not homeomorphic to another.  
    \begin{enumerate}
      \item $M_0 = S^2$ (sphere). $M_1 = T^2$ (torus). $M_2$ is a donut with two holes. $M_3$ has three holes, and so on. 
      \item $N_1$ is the Mobius strip. $N_2$ is the Klein bottle. 
    \end{enumerate}
  \end{example}

\subsection{Retractions}

  \begin{definition}[Retraction]
		\label{def:retraction}
    Given topological space $X$ and $A \subset X$, a \textbf{retraction} $f: X \to A$ is a continuous map such that $f(a) = a$ for all $a \in A$. 
  \end{definition}

  So, the restriction of retractions onto their codomain is the identity map. 

  \begin{theorem}[Retractions are Quotient Maps]
    A retraction $r$ is a quotient map. 
  \end{theorem}
  \begin{proof}
    We know that the canonical injection $A \subset X$ is continuous, and $r \circ \iota = \mathrm{id}_A$. Therefore, by the lemma, $r$ must be a quotient map. 
  \end{proof}

\subsection{Exercises}

  \begin{exercise}[Munkres 18.1]
    Prove that for functions $f : \mathbb{R} \to \mathbb{R}$, the $\epsilon$-$\delta$ definition of continuity implies the open set definition.
  \end{exercise}

  \begin{exercise}[Munkres 18.2]
    Suppose that $f : X \to Y$ is continuous. If $x$ is a limit point of the subset $A$ of $X$, is it necessarily true that $f(x)$ is a limit point of $f(A)$?
  \end{exercise}
  \begin{solution}
    No. Consider $X = Y = \mathbb{R}$ with the Euclidean topology and let $f(x) = 0$. Consider $A = (-1, 1) \implies f(A) = \{0\}$. $1$ is a limit point of $A$ but $f(1) = 0$ is not a limit point of $\{0\}$ since it's an isolated point, i.e. for any punctured open neighborhood $U_0^\circ = U_0 \setminus \{0\}$, 
    \begin{equation}
      U_0^\circ \cap f(A) = (U_0 \setminus \{0\}) \cap \{0\} = \emptyset
    \end{equation}
  \end{solution}

  \begin{exercise}[Munkres 18.3]
    Let $X$ and $X'$ denote a single set in the two topologies $\mathcal{T}$ and $\mathcal{T}'$, respectively. Let $i : X' \to X$ be the identity function.
    \begin{enumerate}
      \item Show that $i$ is continuous $\Leftrightarrow \mathcal{T}'$ is finer than $\mathcal{T}$.
      \item Show that $i$ is a homeomorphism $\Leftrightarrow \mathcal{T}' = \mathcal{T}$.
    \end{enumerate}
  \end{exercise}

  \begin{exercise}[Munkres 18.4]
    Given $x_0 \in X$ and $y_0 \in Y$, show that the maps $f : X \to X \times Y$ and $g : Y \to X \times Y$ defined by
    \begin{align*}
      f(x) = x \times y_0 \quad \text{and} \quad g(y) = x_0 \times y
    \end{align*}
    are imbeddings.
  \end{exercise}

  \begin{exercise}[Munkres 18.5]
    Show that the subspace $(a, b)$ of $\mathbb{R}$ is homeomorphic with $(0, 1)$ and the subspace $[a, b]$ of $\mathbb{R}$ is homeomorphic with $[0, 1]$.
  \end{exercise}

  \begin{exercise}[Munkres 18.6]
    Find a function $f : \mathbb{R} \to \mathbb{R}$ that is continuous at precisely one point.
  \end{exercise}

  \begin{exercise}[Munkres 18.7]
    \begin{enumerate}
      \item Suppose that $f : \mathbb{R} \to \mathbb{R}$ is ``continuous from the right,'' that is,
      \begin{align*}
        \lim_{x \to a^+} f(x) = f(a),
      \end{align*}
      for each $a \in \mathbb{R}$. Show that $f$ is continuous when considered as a function from $\mathbb{R}_\ell$ to $\mathbb{R}$.
      \item Can you conjecture what functions $f : \mathbb{R} \to \mathbb{R}$ are continuous when considered as maps from $\mathbb{R}$ to $\mathbb{R}_\ell$? As maps from $\mathbb{R}_\ell$ to $\mathbb{R}_\ell$? We shall return to this question in Chapter 3.
    \end{enumerate}
  \end{exercise}
  \begin{solution}
    This definition of continuity at a point in analysis means that for all $\epsilon > 0$ there exists a $\delta > 0$ s.t. $x \in [a, a + \delta) \implies |f(x) - f(a)| < \epsilon$. Assuming this definition holds, let $U_{f(a)} \in \T$ be an open neighborhood of $f(a)$ (w.r.t. Euclidean topology). We wish to show that the preimage is open in $\T_\ell$. 

    By definition of open sets in $\mathbb{R}$, there exists an open ball $(f(a) - \epsilon, f(a) + \epsilon) \subset U_{f(a)}$, and from our analytical definition of continuity there exists a $\delta > 0$ s.t. $f([a, a + \delta)) \subset (f(a) - \epsilon, f(a) + \epsilon)$. Therefore taking the preimage we have 
    \begin{equation}
      [a, a+\delta) \subset f^{-1} \big( (f(a) - \epsilon, f(a) + \epsilon ) \big) \subset f^{-1} \big( U_{f(a)} \big)
    \end{equation}
    Since we can construct an open ball (in $\T_\ell$) $[a, a + \delta) \subset f^{-1} (U_{f(a)})$, by definition $f^{-1} (U_{f(a)})$ is open in $\mathbb{R}_\ell$, making $f$ continuous at $a$. Since this holds for all $a$, $f$ is continuous.\footnote{We could have started off with an arbitrary open set $U$ and chosen an arbitrary $a \in f^{-1} (U)$ to get the same result.} 
  \end{solution}

  \begin{exercise}[Munkres 18.8]
    Let $Y$ be an ordered set in the order topology. Let $f, g : X \to Y$ be continuous.
    \begin{enumerate}
      \item Show that the set $\{x \mid f(x) \leq g(x)\}$ is closed in $X$.
      \item Let $h : X \to Y$ be the function
      \begin{align*}
        h(x) = \min\{f(x), g(x)\}.
      \end{align*}
      Show that $h$ is continuous. [Hint: Use the pasting lemma.]
    \end{enumerate}
  \end{exercise}

  \begin{exercise}[Munkres 18.9]
    Let $\{A_\alpha\}$ be a collection of subsets of $X$; let $X = \bigcup_\alpha A_\alpha$. Let $f : X \to Y$; suppose that $f|A_\alpha$ is continuous for each $\alpha$.
    \begin{enumerate}
      \item Show that if the collection $\{A_\alpha\}$ is finite and each set $A_\alpha$ is closed, then $f$ is continuous.
      \item Find an example where the collection $\{A_\alpha\}$ is countable and each $A_\alpha$ is closed, but $f$ is not continuous.
      \item An indexed family of sets $\{A_\alpha\}$ is said to be \textit{locally finite} if each point $x$ of $X$ has a neighborhood that intersects $A_\alpha$ for only finitely many values of $\alpha$. Show that if the family $\{A_\alpha\}$ is locally finite and each $A_\alpha$ is closed, then $f$ is continuous.
    \end{enumerate}
  \end{exercise}

  \begin{exercise}[Munkres 18.10]
    Let $f : A \to B$ and $g : C \to D$ be continuous functions. Let us define a map $f \times g : A \times C \to B \times D$ by the equation
    \begin{equation}
      (f \times g)(a \times c) = f(a) \times g(c).
    \end{equation}
    Show that $f \times g$ is continuous.
  \end{exercise}
  \begin{solution}
    Let $V$ be an open set in $B \times D$. Then $V$ can be expressed as the union of basis elements in the product topology, with the form 
    \begin{equation}
      V = \bigcup U_B \times U_D
    \end{equation}
    where $U_B \in \mathscr{T}_B$ and $U_D \in \mathscr{T}_D$. Now take the preimage. 
    \begin{equation}
      (f \times g)^{-1} (V) = \bigcup (f \times g)^{-1} (U_B \times U_D) = \bigcup f^{-1}(U_B) \times g^{-1} (U_D)
    \end{equation}
    Since $f, g$ are continuous, $f^{-1} (U_B) \in \mathscr{T}_A$ and $g^{-1} (U_D) \in \mathscr{T}_C$. Therefore, $f^{-1} (U_B) \times g^{-1} (U_D)$ is a basis element of the product topology $\mathscr{T}_{A \times C}$, and its arbitrary union is indeed an open set in $A \times C$. Therefore $f \times g$ is continuous.  
  \end{solution}

  \begin{exercise}[Munkres 18.11]
    Let $F : X \times Y \to Z$. We say that $F$ is \emph{continuous in each variable separately} if for each $y_0$ in $Y$, the map $h : X \to Z$ defined by $h(x) = F(x \times y_0)$ is continuous, and for each $x_0$ in $X$, the map $k : Y \to Z$ defined by $k(y) = F(x_0 \times y)$ is continuous. Show that if $F$ is continuous, then $F$ is continuous in each variable separately.
  \end{exercise}
  \begin{solution}
    Let us define $\iota_{y_0}: X \rightarrow X \times Y$ as the canonical injection $\iota_{y_0}(x) = (x, y_0)$. We first show that this is continuous. First choose an open set $V \in \mathscr{T}_{X \times Y}$, which is of the form 
    \begin{equation}
      V = \bigcup U_X \times U_Y
    \end{equation}
    for open sets $U_X \in \mathscr{T}_X, U_Y \in \mathscr{T}_Y$. Taking the preimage 
    \begin{equation}
      \iota_{y_0}^{-1} (V) = \bigcup \iota_{y_0}^{-1} (U_X \times U_Y)
    \end{equation}
    For each term in the union, note that if $y_0 \not\in U_Y$, then the preimage is $\emptyset$, which is open. If $y_0 \in U_Y$, then the preimage is $U_X$ which is open. Therefore the union of such open sets is open. With this, note that 
    \begin{equation}
      h = F \circ \iota_{y_0}
    \end{equation}
    which is a composition of continuous maps and therefore is continuous. The proof for $k$ is identical. 
  \end{solution}

  \begin{exercise}[Munkres 18.12]
    Let $F : \mathbb{R} \times \mathbb{R} \to \mathbb{R}$ be defined by the equation
    \begin{equation}
      F(x \times y) = \begin{cases}
        xy/(x^2 + y^2) & \text{if } x \times y \neq 0 \times 0, \\
        0              & \text{if } x \times y = 0 \times 0.
      \end{cases}
    \end{equation}
    \begin{enumerate}
      \item Show that $F$ is continuous in each variable separately.
      \item Compute the function $g : \mathbb{R} \to \mathbb{R}$ defined by $g(x) = F(x \times x)$.
      \item Show that $F$ is not continuous.
    \end{enumerate}
  \end{exercise}
  \begin{solution}
    \begin{enumerate}
      \item Fix $y = y_0$. Then if $y_0 = 0$, $h(x) = 0$ which is continuous. If $y_0 \neq 0$, then 
      \begin{equation}
        h(x) = F(x \times y_0) = \frac{y_0 x}{x^2 + y_0^2}
      \end{equation}
      which is the quotient of two polynomials, which are continuous, and the denominator never vanishes since $x^2 + y_0^2 \geq y_0^2 > 0$. Similarly, if we fix $x = x_0$, $k(y) = 0$ if $x_0 = 0$ and 
      \begin{equation}
        k(y) = F(x_0 \times y) = \frac{x_0 y}{y^2 + x_0^2}
      \end{equation}
      which is the quotient of two polynomials where the denominator never vanishes. 

      \item We have 
      \begin{equation}
        g(x) = \begin{cases} 
          F(0) = 0 & \text{ if } x = 0 \\
          F(x \times x) = \frac{x^2}{2 x^2} = \frac{1}{2} & \text{ if } x \neq 0
        \end{cases}
      \end{equation}

      \item We see that $g$ above is not continuous since the preimage of the open set $(-0.25, 0.25)$ is $\{0\}$ which is not open. We can write $g = F \circ \iota$, where $\iota(x) = (x, x)$. We claim that $\iota$ is continuous. Take any basis element $U \times V \in \mathscr{T}_{X \times X}$ where $U, V$ are open in $X$. The preimage consists of all points $x$ that are both in $U$ and $V$, i.e. $\iota^{-1} (U \times V) = U \cap V$, which is open.  Therefore $\iota$ is continuous. If $F$ was continuous, then $F \circ \iota$ would be continuous, but $g$ is not continuous, so $F$ must not be continuous. 
    \end{enumerate}
  \end{solution}

  \begin{exercise}[Munkres 18.13]
    Let $A \subset X$; let $f : A \to Y$ be continuous; let $Y$ be Hausdorff. Show that if $f$ may be extended to a continuous function $g : \bar{A} \to Y$, then $g$ is uniquely determined by $f$.
  \end{exercise}
  \begin{solution}
    Let us consider two such extensions $g, h$. If $A = \overline{A}$, then $g = f = h$ and this is unique. If $A \subsetneq \overline{A}$ then there exists $x \in A^\prime \setminus A$. Since $Y$ is Hausdorff, there exists disjoint open neighborhoods $U \ni g(x), V \ni h(x)$. It is the case that $x \in g^{-1}(U) \cap h^{-1} (V)$ open (since $f, h$ are continuous and so their preimages are open). Since $x$ is a limit point of $A$, there exists a $y \in g^{-1} (U) \cap h^{-1} (V) \cap A$ not equal to $x$. Mapping through through $h, g$ again gives $g(y) \in U$, $h(y) \in V$. But since $y \in A$, the two must agree with $f$, and so $f(y) = g(y) = h(y)$, which contradicts that $Y$ is Hausdorff. 
  \end{solution}

  \begin{exercise}[Math 411 Spring 2025, PS4]
    Suppose $X$ and $Y$ are topological spaces, where $Y$ is Hausdorff, and let $f$ and $g$ be continuous functions from $X$ to $Y$. Prove that the set $S = \{x \in X \mid f(x) = g(x)\}$ is closed.
  \end{exercise}
  \begin{solution}
    We equivalently wish to show that $X \setminus S$ is open. For any $x \in (X \setminus S)$, we have $f(x) \neq g(x) \in Y$. Since $Y$ is Hausdorff, there exists open $U_{f(x)} \ni f(x)$ and open $U_{g(x)} \ni g(x)$ s.t. $U_{f(x)} \cap U_{g(x)} = \emptyset$. Therefore, we can take their preimage $f^{-1} (U_{f(x)}), g^{-1} (U_{g(x)})$ which is open in $X$ by continuity of $f, g$. Furthermore we can take their intersection to get another open neighborhood of $x$. 
    \begin{equation}
     V_x = f^{-1} (U_{f(x)}) \cap g^{-1} (U_{g(x)}) 
    \end{equation}
    We claim that $V_x \cap S = \emptyset$. Assuming not, we have some $s \in V_x \cap S$. Since $s \in V_x$, $f(s) \in U_{f(x)}$ and $g(s) \in U_{g(x)}$, but since $s \in S$, $f(x) = g(x)$ and these map to the same point, contradicting the fact that $U_{f(x)}$ and $U_{g(x)}$ are disjoint. Therefore, our claim holds true, which implies that $V_x \subset X \setminus S$. Therefore, we have proved that all points $x \in (X \setminus S)$ is an interior point, and thus $X \setminus S$ is open.\footnote{This can be shown by letting $X \setminus S$ be the union of all $V_x$ for $x \in (X \setminus S)$ which is open.}
  \end{solution}

  \begin{exercise}[Math 411 Spring 2025, PS5]
    Let $X$ be a topological space, and let $f,g: X \to \mathbb{R}$ be continuous maps.
    \begin{enumerate}
      \item Show that the set $\{x \in X \mid f(x) \leq g(x)\}$ is closed in $X$.
      
      \item Show that the function $h: X \to \mathbb{R}$ given by $h(x) = \max\{f(x),g(x)\}$ is continuous.
    \end{enumerate}
    
    \noindent\textbf{Note:} If you prefer, instead of $\mathbb{R}$ you can use an arbitrary ordered set $Y$ with the order topology in place of $Y$, as in \S18 \#8. (We did not cover the order topology in class; see \S14 if you're curious.)
  \end{exercise}
  \begin{solution}
    For any $y \in \mathbb{R}$. The sets $\{x \in X \mid f(x) > y\}$ and $\{x \in X \mid y > g(x)\}$ are open since they are the preimages of $(y, +\infty)$ and $(-\infty, y)$ respectively. Therefore their intersection is open, i.e. 
    \begin{equation}
      U_y = \{x \in X \mid f(x) > y > g(x)\}
    \end{equation} 
    Therefore their union is also open. 
    \begin{equation}
      U = \bigcup_{y \in \mathbb{R}} U_y = \{x \in X \mid f(x) > g(x)\}
    \end{equation}
    where we can always identify a $y = (f(x) + g(x))/2$. Therefore the complement where $f(x) \leq g(x)$ is closed. 
    
    As for the second part, we will denote $U_f = \{x \in X \mid f(x) \leq g(x)\}$ and $U_g = \{x \in X \mid g(x) \geq f(x) \}$. They are both closed, where $U_g$ can be proved closed by the same symmetric argument. Let us take any closed set $S \subset X$. Since $X = U_f \cup U_g$, we can denote 
    \begin{equation}
      S = S \cap X = S \cap (U_f \cup U_g) = (S \cap U_f) \cup (S \cap U_g) = V_f \cup V_g
    \end{equation}
    where $V_f, V_g$ are closed since we proved before that $U_f, U_g$ are closed. The preimage is 
    \begin{align}
      h^{-1} (S) & = h^{-1} (V_f) \cup h^{-1} (V_g) \\
                 & = g^{-1} (V_f) \cup f^{-1} (V_g)
    \end{align}
    which is the union of open sets and therefore open. 
  \end{solution}

